\section{Discussion}

The results support two related claims. First, central governance signals can be ranked under adversarial strategy injection in a minimal sequential commons. The Fishery Commons study shows that monitoring with sanctions improves held-out collapse under both mutation and live LLM injection, thereby establishing the minimal robustness result. The broader signal sweep then shows that stronger top-down interventions are not interchangeable: adaptive quotas are the best overall policy in medium-tier fishery runs, while temporary closures remain a credible fallback under adversarial-heavy partner mixes. A learned-policy PPO validation in the same tier reaches the same conclusion, so that ranking is not confined to the hand-specified threshold family. This distinction is analytically useful because it separates light-touch enforcement from stronger adaptive intervention.

Second, governance should not be treated only as a purely central instrument. The Harvest Commons study addresses the broader question of governance architecture. Earlier experimental stages included none and bottom-up-only conditions, but the decision-critical comparison reported here narrows to hybrid versus top-down-only control. The central empirical finding is that hybrid governance improves local ecological control relative to top-down-only control, although this improvement carries a measurable welfare cost in several cells. Earlier Harvest pilots and lighter invasion runs suggested this pattern, and the stronger search-over-mutations stress test provides the principal support for it. In the high-power Harvest matrix, hybrid ranks first in all 8 decision cells under the failure-then-patch-health-then-welfare rule and improves patch health in every tier-by-partner block, but several rows show a clearer welfare cost than before. A narrower Harvest PPO validation in the medium tier yields concordant evidence, which indicates that the architecture ranking is not confined to the hand-designed Harvest strategy family. The RL evidence remains secondary because it is medium-tier, fixed-governance, and does not constitute an invasion benchmark. An evaluation restricted to a single collapse statistic would obscure this trade-off. These findings support the move from a single global stock to a second substrate with local interactions and communication.

This framing also explains why the two studies belong in one paper. The Fishery Commons study isolates which central signals work in a highly controlled environment. The Harvest Commons study then asks whether governance remains purely central once local coordination, communication, and targeted intervention are allowed to matter, while evaluating that question through a narrower top-down-only versus hybrid stress test. Commons experiments have long suggested that communication and collective choice can improve outcomes even without a single command-and-control institution \citep{ostrom1992covenants,walker2000collective}. The results reported here extend that intuition to adversarially changing agent populations: once strategy turnover is explicit, governance must be evaluated both as a set of central signals and as a broader architecture for decentralized systems.

The implications extend beyond the specific substrates studied here. More broadly, the paper concerns institution design for self-interested multi-agent systems under uncertainty and turnover. The same protocol can serve as a testbed for shared compute budgets, energy systems, congestion control, spectrum allocation, or decentralized markets, where new behaviors keep entering and governance must balance resilience against performance. Recent generative-society work argues for realistic evaluation under diverse partners and conditions \citep{piatti2024cooperate,govsim_repo,mercer2023willowbrook}; the present results add a more specific lesson: governance should be evaluated both at the level of signal choice and at the level of governance architecture, and stronger invaders can tighten without necessarily overturning those conclusions even when the principal architecture comparison is deliberately narrow and decision-focused.
